% ============================================================
% Bin Derangements and the Gate Width Theorem
% ============================================================

\documentclass[12pt]{amsart}

\usepackage{amsmath,amssymb,amsthm}
\usepackage{booktabs}
\usepackage{microtype}
\usepackage{url}

\newtheorem{theorem}{Theorem}
\newtheorem{lemma}[theorem]{Lemma}
\newtheorem{proposition}[theorem]{Proposition}
\newtheorem{corollary}[theorem]{Corollary}
\newtheorem{conjecture}[theorem]{Conjecture}
\theoremstyle{definition}
\newtheorem{definition}[theorem]{Definition}
\theoremstyle{remark}
\newtheorem{remark}[theorem]{Remark}

\title{Bin Derangements and the Gate Width Theorem}

\author{Alexander S.\ Petty}

\email{alexander.petty@gmail.com}
\date{April 2024}
\subjclass[2020]{11A63, 11A07, 11B83}

\begin{document}
\begin{abstract}
The autocorrelation $R(\ell)$ of a prime repetend vanishes at
lag~$\ell$ precisely when the corresponding group element
$b^{\ell} \in (\mathbb{Z}/p\mathbb{Z})^{*}$ deranges the digit-bin
partition: multiplication by $b^{\ell}$ moves every residue to a
different bin. We prove that for any prime $p > b$ (no
primitive-root hypothesis needed), the bin-deranging set is
exactly
$\{-u/(b{-}u) \bmod p : u = 1, \ldots, b{-}1\}$.
In particular, there are exactly $b - 1$ such elements,
independent of the prime. The gate that controls which Ramanujan
lags survive in the phase-filtered autocorrelation has a universal
width determined solely by the base. The key observation is that
the permutation $r \mapsto br \bmod p$ transforms the contiguous
Beatty bins into residue classes modulo~$b$.

For the non-deranging multipliers $g \ne 1$
(those with $C(g) > 0$), the mean collision count satisfies an
exact formula:
\[
\overline{R} = \frac{Q\bigl(b(Q{-}1) + 2R\bigr)}{b(Q{-}1) + R},
\]
where $Q = \lfloor (p{-}1)/b \rfloor$ and $R = (p{-}1) \bmod b$.
For the spectral class $R = 0$, this simplifies to $\overline{R}
= Q$ exactly. For $Q = 1$ (the smallest primes $p > b + 1$ in any
class), the mean is always~$2$.

Together, these results decompose the autocorrelation into two
independent controls: the base determines how many lags produce
zero autocorrelation ($b - 1$), while the prime determines which
lags are blocked and how strongly the surviving lags contribute.
\end{abstract}
\maketitle

% ============================================================
\section{Introduction}

This paper characterizes the zero set of the autocorrelation.

The digit function $\delta(r) = \lfloor br/p \rfloor$, which
assigns to each residue $r \bmod p$ its leading digit in
base~$b$, has been studied in connection with the distribution
of digits of primes~\cite{mauduit-rivat}, exponential sums
over digit constraints~\cite{shparlinski}, and the statistical
properties of decimal expansions~\cite{hardy-wright}. The
function partitions $\{1, \ldots, p{-}1\}$ into $b$ bins
$B_0, \ldots, B_{b-1}$, where
\[
B_d = \{r \in \{1, \ldots, p{-}1\} : \delta(r) = d\}.
\]
These bins form a Beatty partition: the bin sizes are
$\lfloor (d{+}1)p/b \rfloor - \lfloor dp/b \rfloor$,
determined by the three-distance theorem for the sequence
$\{r/p\}$ at scale $1/b$~\cite{sos}.

The autocorrelation of the repetend of $1/p$ at lag~$\ell$
counts positions $j$ where the digit at position~$j$ equals
the digit at position $j + \ell$:
\[
R(\ell) = \bigl|\{j \in \{0, \ldots, L{-}1\} :
\delta(b^j) = \delta(b^{j+\ell})\}\bigr|
= \sum_{d=0}^{b-1} |B_d \cap g B_d|,
\]
where $g = b^{\ell} \bmod p$ and the intersection is under
multiplication in $(\mathbb{Z}/p\mathbb{Z})^{\times}$.

The question of when $R(\ell) = 0$ reduces to a combinatorial
question about the Beatty partition: for which group elements
$g$ does multiplication by $g$ move every residue to a
different bin? Such elements are \emph{derangements} of the
bin partition, in the sense of the classical theory of
permutation derangements~\cite{stanley}. We determine their
exact number and give an explicit formula.

% ============================================================
\section{The Bin Partition}

The bins $B_0, \ldots, B_{b-1}$ are contiguous intervals of
$\{1, \ldots, p{-}1\}$ determined by the floor function
$\delta(r) = \lfloor br/p \rfloor$. Their sizes are controlled
by the division $p - 1 = bQ + R$, where
$Q = \lfloor (p{-}1)/b \rfloor$ and $R = (p{-}1) \bmod b$:
exactly $R$ bins have size $Q + 1$ and $b - R$ bins have
size~$Q$. This distribution is a special case of the
three-distance theorem~\cite{sos} applied to the Farey-like
partition of $[0, 1)$ by the points $r/p$.

When $p \le b + 1$, each bin contains at most one residue,
the partition is injective, and the autocorrelation is trivially
zero at all non-zero lags. When $p > b + 1$, multiple residues
share bins, and non-trivial autocorrelation structure emerges.

% ============================================================
\section{Bin Derangements}

\begin{definition}
An element $g \in (\mathbb{Z}/p\mathbb{Z})^{*}$ is a
\emph{bin-deranging multiplier} (or \emph{bin derangement}) if
\[
g B_d \cap B_d = \varnothing \quad \text{for every } d \in
\{0, \ldots, b{-}1\}.
\]
That is, multiplication by $g$ moves every residue out of its
digit bin. The terminology follows the classical notion of a
derangement (a permutation with no fixed points)~\cite{stanley},
extended to the partition setting: no bin is preserved setwise.
\end{definition}

\begin{proposition}\label{prop:derangement-equiv}
For a prime $p$ where $b$ is a primitive root,
$R(\ell) = 0$ if and only if $g = b^{\ell} \bmod p$ is a
bin-deranging multiplier.
\end{proposition}

\begin{proof}
$R(\ell) = \sum_d |B_d \cap g B_d|$ where $g = b^{\ell}$.
Each term is non-negative, so $R(\ell) = 0$ if and only if
every term vanishes, which is the bin-derangement condition.
\end{proof}

The bin-derangement condition has a geometric interpretation.
The bins $B_0, \ldots, B_{b-1}$ are contiguous intervals in
$\mathbb{Z}/p\mathbb{Z}$ (since $\delta$ is a floor function of
a linear map). Multiplication by $g$ stretches and wraps these
intervals. A bin derangement is a multiplier that stretches every
interval far enough that it no longer overlaps its original
position.

% ============================================================
\section{The Gate Width Theorem}

\begin{lemma}[Conjugation]\label{lem:conjugation}
For each $r \in \{1, \ldots, p{-}1\}$, let $x = [br]_p$ (the
least positive residue of $br$ modulo $p$). Then
$br = p\,\delta(r) + x$, and consequently
\[
\delta(r) = \delta(s)
\quad\Longleftrightarrow\quad
[br]_p \equiv [bs]_p \pmod{b}.
\]
Under the permutation $r \mapsto [br]_p$, the digit bins are
exactly the residue classes modulo~$b$.
\end{lemma}

\begin{proof}
Since $p \nmid br$, the least positive residue of $br$ modulo $p$
is the remainder after division: $br = p\lfloor br/p \rfloor +
[br]_p = p\,\delta(r) + x$. Reducing modulo $b$:
$0 \equiv p\,\delta(r) + x \pmod{b}$. Since $\gcd(p,b) = 1$,
the residue class of $x$ modulo $b$ determines $\delta(r)$
uniquely.
\end{proof}

\begin{lemma}\label{lem:collision-congruence}
For every $g \in \mathbb{F}_p^{\times}$,
\[
C(g) = \#\{x \in \{1, \ldots, p{-}1\} :
x \equiv [gx]_p \pmod{b}\}.
\]
\end{lemma}

\begin{proof}
The map $r \mapsto x = [br]_p$ is a permutation of
$\{1, \ldots, p{-}1\}$. Apply Lemma~\ref{lem:conjugation} to
the pair $r$ and $gr \bmod p$.
\end{proof}

\begin{lemma}\label{lem:no-collision}
Fix $g \ne 1$ and define $c \equiv b(1{-}g)^{-1} \pmod{p}$ with
$1 \le c \le p{-}1$. If $1 \le c \le b{-}1$, then $C(g) = 0$.
\end{lemma}

\begin{proof}
The key geometric idea is that when $c$ is small (between $1$ and
$b-1$), the shifted sequence $mc \bmod p$ stays too close to zero
to reach the reference sequence $mb$, preventing any collision.

A collision gives $y = [gx]_p$ with $y = x + mb$ for some
integer~$m$, and $x \equiv -mc \pmod{p}$.

If $m \ge 0$: from $x + mb \le p{-}1$,
$1 \le x + mc \le p{-}1{-}m(b{-}c) < p$, but $x + mc \equiv 0
\pmod{p}$, a contradiction.

If $m < 0$: write $m = -n$ with $n \ge 1$; from $x - nb \ge 1$,
$0 < x - nc = (x{-}nb) + n(b{-}c) \le p{-}1$,
again contradicting $x - nc \equiv 0 \pmod{p}$.
\end{proof}

\begin{lemma}\label{lem:collision-exists}
If $c \ge b + 1$, then $C(g) \ge 1$.
\end{lemma}

\begin{proof}
Conversely, when $c$ is large enough ($c \ge b+1$), the gap between
the shifted and reference sequences is wide enough that a collision
point can be constructed explicitly.

Take $x = p - c$. Then $1 \le x \le p{-}b{-}1$, so
$y := x + b = p - c + b$ lies in $\{1, \ldots, p{-}1\}$ with
$y \equiv x \pmod{b}$. Since $x \equiv -c \pmod{p}$ and
$c(1{-}g) \equiv b$, we get $(1{-}g)x \equiv -b$, hence
$gx \equiv x + b = y \pmod{p}$. So $x$ contributes a collision.
\end{proof}

\begin{theorem}[Gate width]\label{thm:gate-width}
Let $p$ be prime and $2 \le b < p$. Then
\[
\{g \in \mathbb{F}_p^{\times} : C(g) = 0\}
= \left\{-\frac{u}{\,b - u\,} \pmod{p} :
u = 1, \ldots, b{-}1\right\}.
\]
In particular, there are exactly $b - 1$ bin-deranging elements.
No primitive-root hypothesis is needed.
\end{theorem}

\begin{proof}
By Lemmas~\ref{lem:no-collision} and~\ref{lem:collision-exists},
$C(g) = 0$ if and only if $c = b(1{-}g)^{-1} \in \{1, \ldots,
b{-}1\}$. Solving: $g = 1 - b/c$. Setting $u = b - c$ gives
$g = -u/(b{-}u) \pmod{p}$ for $u = 1, \ldots, b{-}1$. These
are distinct: $-u_1/(b{-}u_1) \equiv -u_2/(b{-}u_2) \pmod{p}$
implies $b(u_1 - u_2) \equiv 0 \pmod{p}$, and since $p > b$
this gives $u_1 = u_2$.
\end{proof}

\begin{corollary}\label{cor:ratio-set}
$\displaystyle
\left|\bigcup_{d=0}^{b-1} B_d / B_d\right| = p - b$.
\end{corollary}

\begin{remark}
The key step is the conjugation
(Lemma~\ref{lem:conjugation}): the permutation $r \mapsto [br]_p$
transforms the contiguous Beatty bins into residue classes
modulo~$b$. The partition geometry linearizes completely under
this change of variables. The squeezing argument in
Lemmas~\ref{lem:no-collision} and~\ref{lem:collision-exists}
then determines the zero set exactly.
\end{remark}

\begin{remark}
The deranging multipliers $g = -u/(b{-}u) \pmod{p}$ form a
rational family. For $u = b/2$ (when $b$ is even),
$g \equiv -1 \pmod{p}$, the complement element. The remaining
derangers come in inverse pairs: $-u/(b{-}u)$ and $-(b{-}u)/u$
are mutual inverses modulo~$p$.
\end{remark}

\begin{remark}
The number $b - 1$ depends only on the base. It does not depend
on the prime, the spectral class, or the scale. The gate width
is a property of the coordinate system, not of the object being
observed. What varies with the prime is \emph{which} $b - 1$
lags are blocked, and \emph{how strongly} the surviving lags
contribute.
\end{remark}

% ============================================================
\section{The Sum Rule and Mean Formula}

The autocorrelation satisfies a sum rule analogous to
Parseval's identity for Ramanujan sums~\cite{ramanujan},
connecting the collision count to the spectral power at
frequency zero.

\begin{proposition}[Sum rule]\label{prop:sum-rule}
\[
\sum_{g \in \mathbb{F}_p^{\times}} C(g)
= \Phi(0) = \sum_{d=0}^{b-1} n_d^2,
\]
where $n_d = |B_d|$ is the size of the $d$-th digit bin.
\end{proposition}

\begin{proof}
$\sum_g C(g) = \sum_g \sum_r [\delta(r) = \delta(gr)]
= \sum_r \sum_g [\delta(gr) = \delta(r)]$.
For each $r$, as $g$ ranges over $\mathbb{F}_p^{\times}$, the
element $gr$ ranges over all of $\mathbb{F}_p^{\times}$.
The inner sum counts how many residues share a bin with $r$,
which is $n_{\delta(r)}$. Summing over $r$:
$\sum_r n_{\delta(r)} = \sum_d n_d^2$.
\end{proof}

The bin sizes for a prime $p > b + 1$ are determined by the division
$p - 1 = bQ + R$, where $Q = \lfloor (p{-}1)/b \rfloor$ and $R
= (p{-}1) \bmod b$. There are $R$ bins of size $Q + 1$ and $b -
R$ bins of size $Q$. The spectral power at zero is therefore
\[
\Phi(0) = R(Q+1)^2 + (b-R)Q^2 = bQ^2 + R(2Q + 1).
\]

\begin{theorem}[Mean constructive value]\label{thm:mean}
The mean of $C(g)$ over the $p - b - 1$ constructive
multipliers (those $g \ne 1$ with $C(g) > 0$) is
\[
\overline{C} = \frac{Q\bigl(b(Q{-}1) + 2R\bigr)}{b(Q{-}1) + R}.
\]
When $b$ is a primitive root modulo $p$, this equals the mean of
$R(\ell)$ over constructive lags.
\end{theorem}

\begin{proof}
There are $p - 2$ nonzero multipliers, of which $b - 1$ have
$C(g) = 0$ by Theorem~\ref{thm:gate-width}. The number of
constructive multipliers is $p - 2 - (b-1) = p - b - 1
= b(Q{-}1) + R$.

The sum over constructive multipliers equals
$\Phi(0) - C(1) = bQ^2 + R(2Q+1) - (p-1) = bQ^2 - bQ + 2RQ
= Q(b(Q{-}1) + 2R)$.

Dividing: $\overline{C} = Q(b(Q{-}1) + 2R) / (b(Q{-}1) + R)$.
\end{proof}

\begin{corollary}\label{cor:special-cases}
\begin{enumerate}
\item For $R = 0$ (class $p \equiv 1 \pmod{b}$):
      $\overline{C} = Q$ exactly.
\item For $Q = 1$ (the smallest primes $p > b + 1$):
      $\overline{C} = 2$.
\item As $Q \to \infty$ (large primes in any class):
      $\overline{C} \to Q \sim p/b$.
\end{enumerate}
\end{corollary}

\begin{proof}
(1) Set $R = 0$: the formula gives $Q \cdot b(Q{-}1) / b(Q{-}1)
= Q$.
(2) Set $Q = 1$: the numerator is $1 \cdot (0 + 2R) = 2R$ and
the denominator is $0 + R = R$, giving $2$.
(3) For large $Q$: $Q(bQ + 2R) / (bQ + R) \to Q$ since $R < b$
is bounded.
\end{proof}

% ============================================================
\section{Computational Verification}

All results have been verified using the \texttt{nfield} software
engine across multiple bases.

\begin{table}[h]
\centering
\begin{tabular}{rrrrrrrr}
\toprule
$b$ & $p$ & $r$ & $Q$ & $R$ & zeros & $\overline{R}$ actual
& $\overline{R}$ formula \\
\midrule
$10$ & $17$ & $7$ & $1$ & $6$ & $9$ & $2.00$ & $2.00$ \\
$10$ & $61$ & $1$ & $6$ & $0$ & $9$ & $6.00$ & $6.00$ \\
$10$ & $97$ & $7$ & $9$ & $6$ & $9$ & $9.63$ & $9.63$ \\
$10$ & $131$ & $1$ & $13$ & $0$ & $9$ & $13.00$ & $13.00$ \\
$10$ & $179$ & $9$ & $17$ & $8$ & $9$ & $17.81$ & $17.81$ \\
$10$ & $193$ & $3$ & $19$ & $2$ & $9$ & $19.21$ & $19.21$ \\
$6$ & $11$ & $5$ & $1$ & $4$ & $5$ & $2.00$ & $2.00$ \\
$6$ & $41$ & $5$ & $6$ & $4$ & $5$ & $6.24$ & $6.24$ \\
$6$ & $89$ & $5$ & $14$ & $4$ & $5$ & $14.11$ & $14.11$ \\
$7$ & $17$ & $3$ & $2$ & $2$ & $6$ & $2.50$ & $2.50$ \\
$7$ & $41$ & $6$ & $5$ & $5$ & $6$ & $5.42$ & $5.42$ \\
$7$ & $97$ & $6$ & $13$ & $5$ & $6$ & $13.11$ & $13.11$ \\
$12$ & $17$ & $5$ & $1$ & $4$ & $11$ & $2.00$ & $2.00$ \\
$12$ & $67$ & $7$ & $5$ & $6$ & $11$ & $5.17$ & $5.17$ \\
\bottomrule
\end{tabular}
\caption{Gate width and mean formula verified across bases. The
zero count equals $b - 1$ in every case, and the mean formula is
exact.}
\label{tab:verification}
\end{table}

% ============================================================
\section{The Two Controls}

The autocorrelation of a prime repetend is governed by two
independent parameters:

\begin{enumerate}
\item \textbf{The base} determines the gate width: exactly
$b - 1$ lags are blocked, regardless of which prime is being
analyzed. This is a property of the coordinate system (the
positional representation), not of the prime.

\item \textbf{The prime} determines the gate pattern and
intensity. Among the $p - b - 1$ constructive multipliers,
the distribution of $C(g)$ values carries the structural
fingerprint of the prime. The mean intensity scales as
$\overline{C} \approx Q \approx p/b$ for large primes, but the
distribution around the mean encodes the bin geometry.
\end{enumerate}

This decomposition separates two independent controls.
The number of blocked lags is universal ($b - 1$), depending
only on the base. The positions of blocked lags depend on~$p$
(not just its residue class modulo~$b$), and the mean
intensity at constructive lags has an exact closed form
in~$(Q, R)$.

% ============================================================
\section{Remarks}

\subsection*{The linearization}

The proof reveals that the partition geometry, which resisted
Fourier, combinatorial, and representation-theoretic approaches,
linearizes completely under the change of variables $r \mapsto
[br]_p$. In $r$-space, the bins are contiguous Beatty intervals.
In $x = [br]_p$ space, the same bins are the residue classes
modulo~$b$. The collision condition $\delta(r) = \delta(gr)$
becomes the congruence $x \equiv [gx]_p \pmod{b}$, and the
squeezing argument in Steps~3 and~4 completes the proof.

\subsection*{The explicit deranging set}

The deranging multipliers $g = -u/(b{-}u) \pmod{p}$ for
$u = 1, \ldots, b{-}1$ form a rational family. For $u = b/2$
(when $b$ is even), $g \equiv -1 \pmod{p}$, the complement
element. The remaining derangers come in inverse pairs:
$-u/(b{-}u)$ and $-(b{-}u)/u$ are mutual inverses modulo~$p$.

\subsection*{Generality}

The theorem holds for all primes $p > b$, not only those where
$b$ is a primitive root. The original computational verification
used the primitive-root hypothesis because the autocorrelation
$R(\ell)$ is defined on the orbit of $b$. The algebraic proof
bypasses the orbit entirely, working directly with the digit
function and the collision count $C(g)$.

% ============================================================
\begin{thebibliography}{99}

\bibitem{hardy-wright}
G.~H.~Hardy and E.~M.~Wright,
\emph{An Introduction to the Theory of Numbers},
6th~ed., Oxford University Press, 2008.

\bibitem{mauduit-rivat}
C.~Mauduit and J.~Rivat,
Sur un probl\`eme de Gelfond: la somme des chiffres des
nombres premiers,
\emph{Ann.\ of Math.}\ \textbf{171} (2010), 1591--1646.

\bibitem{ramanujan}
S.~Ramanujan,
On certain trigonometrical sums and their applications in the
theory of numbers,
\emph{Trans.\ Cambridge Philos.\ Soc.}\ \textbf{22} (1918),
259--276.

\bibitem{shparlinski}
I.~E.~Shparlinski,
Character sums over shifted primes,
\emph{Mathematika}\ \textbf{56} (2010), 363--372.

\bibitem{sos}
V.~T.~S\'os,
On the distribution mod 1 of the sequence $n\alpha$,
\emph{Ann.\ Univ.\ Sci.\ Budapest.\ E\"otv\"os Sect.\ Math.}\
\textbf{1} (1958), 127--134.

\bibitem{stanley}
R.~P.~Stanley,
\emph{Enumerative Combinatorics}, vol.~1,
2nd~ed., Cambridge University Press, 2012.

\bibitem{nfield}
A.~S.~Petty,
\emph{nfield: A structural analysis engine for fractional field
invariants},
software.
\url{https://github.com/alexspetty/nfield}

\end{thebibliography}

\end{document}
